\chapter{The Pseudo++ Toolchain}
\label{chap:implementation}

This chapter describes the design and implementation of the Pseudo++ toolchain,
a complete software pipeline for writing, executing, debugging, and transpiling
Romanian-style pseudocode as taught at the Romanian high-school level.
The toolchain is written in C (core library) and JavaScript (web editor,
Tree-sitter grammar), and is distributed as a native command-line binary
and a browser-based editor that comes bundled with the core implementation
compiled to WebAssembly.

% ─────────────────────────────────────────────────────────────────────────────
\section{System Architecture}
\label{sec:architecture}

The Pseudo++ toolchain is structured as four layers of independent modules,
each with a narrow public interface. Figure~\ref{fig:modules} shows the
compile-time dependency graph of all modules grouped by layer.

\begin{figure}[h]
\centering
\resizebox{0.85\linewidth}{!}{%
\begin{tikzpicture}[
    product/.style = {
        draw, line width=1.2pt, rounded corners=6pt,
        font=\small, align=center, inner sep=10pt
    },
    build/.style = {
        draw, line width=1pt, rounded corners=5pt,
        font=\small, align=center, inner sep=8pt
    },
    cmod/.style = {
        draw=gray!45, rounded corners=3pt,
        minimum width=3.2cm, minimum height=0.9cm,
        font=\small\ttfamily, align=center, fill=white
    },
    arr/.style = {-Stealth, line width=1.2pt},
]

%% ── C Core Library ──────────────────────────────────────────────────────
\begin{pgfonlayer}{background}
  \fill[teal!7, rounded corners=10pt]
    (-1.0,-0.65) rectangle (15.2, 5.5);
  \draw[teal!35, line width=1pt, rounded corners=10pt]
    (-1.0,-0.65) rectangle (15.2, 5.5);
\end{pgfonlayer}
\node[font=\small\bfseries\color{teal!55!black}, anchor=north west]
  at (-0.8, 5.45) {C Core Library};

%% Capabilities row (y = 3.9)  — teal tint marks the top layer
\node[cmod, fill=teal!5] (rt) at (1.5,  3.9) {runtime + debugger};
\node[cmod, fill=teal!5] (tr) at (5.2,  3.9) {transpiler};
\node[cmod, fill=teal!5] (chk) at (8.9, 3.9) {check};
\node[cmod, fill=teal!5] (eq) at (12.6, 3.9) {equivalence};

%% Foundation row (y = 2.3)
\node[cmod] (parse) at (1.5,  2.3) {parser (Tree-sitter)};
\node[cmod] (lint)  at (7.0,  2.3) {formatter};
\node[cmod] (io)    at (12.6, 2.3) {I/O abstraction};

%% Utilities row (y = 0.7)
\node[cmod] (env)  at (1.5,  0.7) {environment};
\node[cmod] (str)  at (7.0,  0.7) {string};
\node[cmod] (hmap) at (12.6, 0.7) {hashmap};

%% ── WebAssembly build ────────────────────────────────────────────────
\node[build, fill=violet!8, draw=violet!45,
      minimum width=5.0cm, minimum height=1.0cm]
  (wasm) at (2.9, 7.15)
  {{\bfseries WebAssembly Module}\\[2pt]
   {\footnotesize compiled with Emscripten}};

%% ── Products ─────────────────────────────────────────────────────────
\node[product, fill=blue!7, draw=blue!45,
      minimum width=5.2cm, minimum height=1.55cm]
  (web) at (2.9, 9.4)
  {{\bfseries Web Editor}\\[4pt]
   {\footnotesize Monaco $\cdot$ Debugger $\cdot$ Transpile $\cdot$ Equivalence}};

\node[product, fill=orange!9, draw=orange!50,
      minimum width=4.4cm, minimum height=1.55cm]
  (cli) at (12.0, 9.4)
  {{\bfseries CLI}\\[4pt]
   {\footnotesize run $\cdot$ format $\cdot$ check $\cdot$ parse}\\[1pt]
   {\footnotesize debug $\cdot$ transpile $\cdot$ equivalence}};

%% ── Arrows ───────────────────────────────────────────────────────────
\draw[arr, blue!55]
  (web.south) -- node[right, font=\footnotesize\itshape, text=blue!55]
  {JavaScript calls} (wasm.north);

\draw[arr, violet!60]
  (wasm.south) -- node[right, font=\footnotesize\itshape, text=violet!55]
  {native C API} (2.9, 5.5);

\draw[arr, orange!70!black]
  (cli.south) -- (12.0, 5.5);

\end{tikzpicture}%
}
\caption{Compile-time dependency graph of the Pseudo++ modules, grouped by layer.
  The C core library (teal) is compiled to both a native CLI binary and a WebAssembly
  module loaded by the web editor.}
\label{fig:modules}
\end{figure}

The modules and their responsibilities are:

\begin{description}
  \item[\texttt{formatter}] Normalises the raw source text before parsing: it
        replaces Unicode mathematical symbols (e.g.\ $\leq$, $\neq$) with
        their ASCII equivalents and re-indents the code structurally.

  \item[\texttt{check}] Performs static semantic analysis on the CST after
        parsing, reporting diagnostics (errors and warnings with codes such as
        \texttt{W001}) for issues like unused variables or suspicious patterns.

  \item[\texttt{parser}] Wraps the Tree-sitter incremental parser with the
        custom \texttt{pseudo} grammar, exposes the resulting concrete syntax
        tree (CST), and provides helper utilities for navigating and extracting
        text from CST nodes.

  \item[\texttt{runtime}] Implements the tree-walking interpreter.  It
        subdivides into three sub-modules: \texttt{value} (runtime values and
        arithmetic), \texttt{environment} (variable store), and
        \texttt{interpreter} (execution engine with explicit stack).

  \item[\texttt{debugger}] Extends the interpreter with variable inspection
        and JSON serialisation of runtime state for the browser front end.

  \item[\texttt{transpiler}] Converts the CST to equivalent C, C++, or Pascal
        source, performing a two-pass analysis: first collecting variable types,
        then emitting the target code.

  \item[\texttt{equivalence}] Implements loop-equivalence transformations,
        converting a loop construct of one kind (e.g.\ \texttt{for}) into any
        of the other three (\texttt{while}, \texttt{do\_while}, \texttt{repeat})
        while preserving semantics.

  \item[\texttt{io}] Provides an abstract I/O interface with two concrete
        backends: \texttt{io\_stdio} for the CLI and \texttt{io\_buffered} for
        the WebAssembly build.

  \item[\texttt{wasm}] An Emscripten bridge that exports the runtime,
        transpiler, and equivalence module as \texttt{EMSCRIPTEN\_KEEPALIVE}
        C functions callable from JavaScript.

  \item[\texttt{cli}] The command-line entry point, supporting \texttt{run},
        \texttt{format}, \texttt{check}, \texttt{parse}, \texttt{debug},
        \texttt{transpile}, and \texttt{equivalence} sub-commands.
\end{description}

% ─────────────────────────────────────────────────────────────────────────────
\section{Language Grammar}
\label{sec:grammar}

The grammar for the Pseudo++ language is specified as a Tree-sitter
\cite{treesitter} grammar in JavaScript (\texttt{grammar.js}).  It defines
the notation established by convention in Romanian high-school CS textbooks
and baccalaureate examination papers \cite{romanian-curriculum}.

\subsection{Statement Forms}

The top-level production \texttt{program} is a repetition of \texttt{stmt}
nodes.  Each statement is one of:

\begin{itemize}
  \item \textbf{Assignment} (\texttt{assign}):
        \texttt{x <- expr} --- assigns the value of \textit{expr} to variable
        \texttt{x}.
  \item \textbf{Swap} (\texttt{swap}):
        \texttt{x <-> y} --- exchanges the values of \texttt{x} and \texttt{y}.
  \item \textbf{Read} (\texttt{read}):
        \texttt{citeste x, y, z} --- reads values from standard input.
  \item \textbf{Write} (\texttt{write}):
        \texttt{scrie expr1, expr2} --- writes values to standard output.
  \item \textbf{Conditional} (\texttt{if}):
        \texttt{daca} \textit{cond} \texttt{atunci} \ldots\
        [\texttt{altfel} \ldots] \texttt{sf}.
  \item \textbf{Counting loop} (\texttt{for}):
        \texttt{pentru i <- start, end[, step] executa} \ldots\ \texttt{sf}.
  \item \textbf{Pre-condition loop} (\texttt{while}):
        \texttt{cat timp} \textit{cond} \texttt{executa} \ldots\ \texttt{sf}.
  \item \textbf{Post-condition loop} (\texttt{do\_while}):
        \texttt{executa} \ldots\ \texttt{cat timp} \textit{cond}.
  \item \textbf{Repeat-until loop} (\texttt{repeat}):
        \texttt{repeta} \ldots\ \texttt{pana cand} \textit{cond}.
  \item \textbf{Compound statement} (\texttt{multi\_stmt}): two or more simple
        statements separated by semicolons, used where only a single statement
        is syntactically permitted.
\end{itemize}

\subsection{Expression Hierarchy}

Expressions follow a standard precedence hierarchy encoded via Tree-sitter's
\texttt{prec.left} operator:

\begin{center}
\begin{tabular}{c l l}
\hline
\textbf{Level} & \textbf{Operator(s)} & \textbf{Node type} \\
\hline
1 (lowest) & \texttt{SAU / sau}                        & \texttt{or\_expr} \\
2          & \texttt{SI / si}                           & \texttt{and\_expr} \\
3          & \texttt{NOT / not}                         & \texttt{not\_expr} \\
4          & \texttt{= != < <= > >=}                    & \texttt{compare\_expr} \\
5          & \texttt{+ -}                               & \texttt{add\_expr} \\
6          & \texttt{* / \%}                            & \texttt{mul\_expr} \\
7 (highest)& \texttt{-(x)}, \texttt{[x]}, $\sqrt{x}$ & \texttt{neg\_expr}, \texttt{floor}, \texttt{sqrt\_expr} \\
\hline
\end{tabular}
\end{center}

Atoms (\texttt{atom}) are integer literals, floating-point literals, string
literals, or identifiers.  Parenthesised sub-expressions (\texttt{paren}) reset
precedence.  The floor operator uses bracket notation (\texttt{[expr]}), which
maps to integer truncation, while the square-root operator uses the radical
symbol ($\sqrt{\phantom{x}}$).

% ─────────────────────────────────────────────────────────────────────────────
\section{Source Normalisation}
\label{sec:formatter}

The Romanian high-school pseudocode notation, established by convention
through textbooks and baccalaureate examination papers \cite{romanian-curriculum},
uses a rich set
of Unicode characters: relational symbols ($\leq$, $\neq$, $\geq$), assignment and
swap arrows ($\leftarrow$, $\leftrightarrow$), box-drawing characters for
structural indentation (box-drawing bar), a filled square ($\blacksquare$) as the
block-closing symbol instead of the keyword \texttt{sf}, and letters with
Romanian diacritics.  This notation is the \emph{canonical} form of the
language --- it is not incidental variation but the deliberate typographic
standard taught in schools.

Formalising this notation directly into a parser grammar would tightly couple
the grammar to a specific Unicode encoding and make the language impractical to
type in a plain-text environment.  The toolchain therefore introduces a
deliberate split: the \emph{source language} that users write (and that the
parser accepts) uses ASCII-compatible tokens, while the \emph{display language}
visible in the web editor renders those tokens back into the canonical Unicode
form via ligatures and glyph substitutions.  This rendering is purely visual
and optional: a toggle button in the editor enables or disables it, leaving
the underlying source file unchanged in either case.

The formatter is the bridge between the two representations.  It accepts source in
either form --- canonical Unicode or ASCII --- and produces the normalised ASCII
form consumed by the parser.  The normalisation pipeline consists of two passes:

\paragraph{Pass 1 --- Symbol Substitution.}
The formatter maintains a static hash map that maps Unicode byte sequences to their
ASCII equivalents (Table~\ref{tab:substitutions}).  The pass iterates through
the source byte-by-byte, scanning the substitution map at each position for the
longest matching key (a greedy, maximum-munch strategy).  When a match is found
the corresponding replacement string is appended to the output and the cursor
advances by the match length; otherwise the current byte is copied verbatim.

\begin{table}[h]
\centering
\small
\begin{tabular}{l l l}
\hline
\textbf{Canonical input} & \textbf{Category} & \textbf{ASCII output} \\
\hline
$\leq$, $\neq$, $\geq$        & Relational symbols   & \texttt{<=}, \texttt{!=}, \texttt{>=} \\
$\leftarrow$                   & Assignment arrow     & \texttt{<-} \\
$\leftrightarrow$              & Swap arrow           & \texttt{<->} \\
$\blacksquare$ (U+25A0)        & Block-close square   & \texttt{sf} \\
box-drawing bar (U+2502)       & Structural indent    & 4 spaces \\
\texttt{\textquotedblleft}, \texttt{\quotedblbase} & Curly quotes & \texttt{"} \\
\textit{ă}, \textit{â}, \textit{î}, \textit{ș}, \textit{ț} & Romanian diacritics & \texttt{a}, \texttt{a}, \texttt{i}, \texttt{s}, \texttt{t} \\
\hline
\end{tabular}
\caption{Selected substitution rules applied by the formatter in Pass 1.
  The canonical Unicode notation is the standard form used in Romanian
  high-school textbooks and exam papers; the ASCII output is the form
  accepted by the Tree-sitter grammar.}
\label{tab:substitutions}
\end{table}

\paragraph{Pass 2 --- Structural Re-indentation.}
After symbol substitution the source may have inconsistent indentation,
particularly when the box-drawing indentation markers have been replaced by
spaces.  The second pass strips all leading whitespace from every line and
recomputes indentation from scratch based on a set of \emph{indent triggers}:

\begin{itemize}
  \item \textit{Indent-after} lines --- lines whose last word is \texttt{atunci},
        \texttt{executa}, \texttt{altfel}, or whose first word is \texttt{repeta}
        --- increase the indentation level for the next line.
  \item \textit{Dedent-before} lines --- \texttt{sf} alone, \texttt{altfel} as
        the first word, and the closing \texttt{pana cand} of a repeat loop ---
        decrease the indentation level before the line is emitted.
\end{itemize}

The result is a clean, consistently-indented ASCII source that the Tree-sitter
parser can analyse reliably.

% ─────────────────────────────────────────────────────────────────────────────
\section{Parsing and CST Construction}
\label{sec:parsing-ast}

The \texttt{parser} module wraps the Tree-sitter C API.  Its public interface
centres on an opaque \texttt{parser\_t} struct that owns a
\texttt{TSParser} (the Tree-sitter parser instance), a \texttt{TSTree} (the
last parse result), and a copy of the normalised source string.

\paragraph{Parsing.}
\texttt{parser\_parse()} first copies the source, then calls
\texttt{ts\_parser\_parse\_string()}.  Tree-sitter guarantees a result tree
even on malformed input, using \texttt{ERROR} and \texttt{MISSING} synthetic
nodes to indicate problem sites.  After the C API call, the function performs a
depth-first walk of the resulting tree looking for the first such error node;
if one is found, it reports the row and column of the error and constructs a
human-readable diagnostic message.

\paragraph{Node Helpers.}
Because the interpreter and transpiler perform extensive tree traversal,
\texttt{parser.c} also provides a set of helper functions used throughout the
codebase:

\begin{itemize}
  \item \texttt{parser\_node\_is\_type(node, type)} --- string-compares the
        node's grammar type against a constant.
  \item \texttt{parser\_child\_by\_field(node, field)} --- wraps
        \texttt{ts\_node\_child\_by\_field\_name}, used to retrieve named
        children such as \texttt{condition}, \texttt{left}, \texttt{right},
        \texttt{var}, \texttt{start}, \texttt{end}, and \texttt{step}.
  \item \texttt{parser\_node\_text(parser, node)} --- extracts the source
        substring corresponding to a node by using its start and end byte
        offsets.
  \item \texttt{parser\_pretty\_tree()} / \texttt{parser\_debug\_tree()} ---
        produce human-readable renderings of the CST for diagnostic purposes.
\end{itemize}

% ─────────────────────────────────────────────────────────────────────────────
\section{Interpreter}
\label{sec:interpreter}

\subsection{Value Representation}
\label{subsec:values}

Runtime values are represented by a heap-allocated, tagged-union type
\texttt{value\_t}:

\begin{lstlisting}[language=C, basicstyle=\ttfamily\small]
struct value {
    value_type_t type;        /* VALUE_INT | VALUE_FLOAT | VALUE_STRING */
    union {
        int64_t   int_val;
        double    float_val;
        string_t* string_val;
    };
};
\end{lstlisting}

\noindent
The three types cover the full range of values expressible in Romanian
baccalaureate pseudocode.  The following design decisions govern arithmetic:

\begin{itemize}
  \item \textbf{Type promotion.}  When one operand is \texttt{VALUE\_FLOAT}
        and the other is \texttt{VALUE\_INT}, the integer is promoted via
        \texttt{value\_to\_float()} before the operation is applied.
  \item \textbf{Division always yields a float.}
        The \texttt{/} operator always returns a \texttt{double}, mirroring
        the mathematical convention used in Romanian textbooks.  Integer
        division is expressed explicitly with the floor operator \texttt{[a/b]}.
  \item \textbf{Modulo operates on integers.}
        The \texttt{\%} operator truncates both operands to \texttt{int64\_t}
        before computing the remainder.
  \item \textbf{Square root returns an integer when the result is exact.}
        If \texttt{sqrt(d) == floor(sqrt(d))}, the result is returned as
        \texttt{VALUE\_INT}.
  \item \textbf{String concatenation.}
        The \texttt{+} operator is overloaded: when both operands are strings
        it performs concatenation; mixed string/numeric operands are a
        type error.
\end{itemize}

Logical operations (\texttt{SI}, \texttt{SAU}, \texttt{NOT}) operate on the
\emph{truthiness} of any value: a numeric zero and an empty string are false;
everything else is true.  They always return \texttt{VALUE\_INT} (1 or 0).

All arithmetic functions follow the same error-reporting convention: an
optional \texttt{value\_error\_t* err} output parameter is set to
\texttt{VALUE\_OK}, \texttt{VALUE\_ERR\_TYPE}, \texttt{VALUE\_ERR\_DIV\_ZERO},
or \texttt{VALUE\_ERR\_NEGATIVE\_SQRT}.

\subsection{Execution Environment}
\label{subsec:environment}

Variables are stored in a flat \texttt{environment\_t}, implemented as an
open-addressing hash table with linear probing.  The table uses three slot
states (\texttt{SLOT\_EMPTY}, \texttt{SLOT\_OCCUPIED}, \texttt{SLOT\_DELETED})
to support deletion without disrupting existing probe chains.  The FNV-1a hash
function \cite{fnv-hash} is used for string keys:

\begin{lstlisting}[language=C, basicstyle=\ttfamily\small]
static uint64_t hash_string(const char* str) {
    uint64_t hash = 14695981039346656037ULL;
    while (*str) {
        hash ^= (uint64_t)(unsigned char)(*str++);
        hash *= 1099511628211ULL;
    }
    return hash;
}
\end{lstlisting}

The table starts with capacity 16.  When the combined occupancy of
\texttt{OCCUPIED} and \texttt{DELETED} slots exceeds a 70\,\% load factor,
\texttt{env\_resize()} doubles the capacity and rehashes all live entries.
Ownership of values is transferred on set: the table takes ownership of
the \texttt{value\_t*} passed to \texttt{env\_set()}, and the previous value
(if any) is destroyed in place.

An important semantic detail is that Pseudo++ variables are \emph{implicitly
initialised}: reading an undeclared variable returns an integer zero and also
creates the binding, matching the convention in Romanian textbook examples
where variables are used before being formally assigned.

\subsection{Expression Evaluation}
\label{subsec:expressions}

Expressions are evaluated by a recursive \texttt{eval\_expr()} function that
walks the CST nodes produced by Tree-sitter.  For binary operators
(\texttt{add\_expr}, \texttt{mul\_expr}, \texttt{compare\_expr}) the function
extracts the named children \texttt{left}, \texttt{op}, and \texttt{right},
evaluates the sub-expressions recursively, reads the operator text, and
dispatches to the appropriate \texttt{value\_*} arithmetic function.

Short-circuit evaluation is implemented for the logical connectives:
\texttt{or\_expr} returns 1 immediately if the left operand is truthy without
evaluating the right; \texttt{and\_expr} returns 0 immediately if the left
operand is falsy.

After each binary expression evaluation the runtime checks whether an error
flag was set; if so, it transitions to the \texttt{EXEC\_ERROR} state and
stores the error message, halting further execution.

\subsection{Control Flow}
\label{subsec:control-flow}

The interpreter uses an \emph{explicit execution stack} rather than recursion.
This design choice serves two goals: it prevents stack overflow for deeply
nested pseudocode, and it allows the debugger to capture and restore the full
execution state as a data structure.

The stack holds up to \texttt{MAX\_STACK\_DEPTH} frames.  Each frame
(\texttt{exec\_frame\_t}) records:

\begin{itemize}
  \item the \texttt{frame\_type\_t} (one of
        \texttt{FRAME\_PROGRAM}, \texttt{FRAME\_BLOCK}, \texttt{FRAME\_IF},
        \texttt{FRAME\_FOR}, \texttt{FRAME\_WHILE}, \texttt{FRAME\_DO\_WHILE},
        \texttt{FRAME\_REPEAT});
  \item the corresponding CST node;
  \item a \texttt{phase} integer and \texttt{child\_idx} counter tracking
        which part of the construct has already been executed;
  \item loop-specific fields (\texttt{loop\_current}, \texttt{loop\_end},
        \texttt{loop\_step}, \texttt{loop\_var}) for the \texttt{for} frame.
\end{itemize}

A compile-time dispatch table maps each frame type to its step function:

\begin{lstlisting}[language=C, basicstyle=\ttfamily\small]
static const frame_step_fn k_step_fns[] = {
    [FRAME_PROGRAM]  = step_program,
    [FRAME_IF]       = step_if,
    [FRAME_FOR]      = step_for,
    [FRAME_WHILE]    = step_while,
    [FRAME_DO_WHILE] = step_do_while,
    [FRAME_REPEAT]   = step_repeat,
    [FRAME_BLOCK]    = step_block,
};
\end{lstlisting}

Each step function advances the frame by one \emph{phase} and returns a
Boolean indicating whether a \emph{visible action} occurred.  Visible actions
include executing a simple statement, evaluating a loop condition, and
terminating the program; invisible actions are internal transitions such as
resetting \texttt{child\_idx} at the start of a new loop iteration.

The public API exposes two execution modes:

\begin{itemize}
  \item \texttt{runtime\_step()} --- advances to the next visible action.
        Used by the step-by-step debugger.
  \item \texttt{runtime\_run()} --- runs to completion or the next input
        request via a tight loop with no overhead.  Used by both the CLI
        \texttt{run} command and the web editor's Run button.
\end{itemize}

The \texttt{EXEC\_NEEDS\_INPUT} state suspends execution when a \texttt{read}
statement encounters an unavailable token.  The web front end supplies the
missing token through \texttt{io\_buffered} and calls \texttt{runtime\_resume()}
to continue.

% ─────────────────────────────────────────────────────────────────────────────
\section{Debugger}
\label{sec:debugger}

The debugger extends the interpreter with step-by-step execution and runtime
state inspection.  The web editor exposes two operations:
\textbf{Step Into} (advance one visible action) and \textbf{Continue} (run to
completion or the next input request).  After each step, the front end displays
the currently executing line, the full variable table, and --- when the step
involved evaluating a condition --- the condition text and its boolean result.

\paragraph{State serialisation.}
The function \texttt{runtime\_get\_variables\_json()} serialises the current
variable table to a JSON array of \texttt{\{name, value, type\}} objects, with
proper JSON string escaping.  This JSON string is the primary mechanism by
which the WebAssembly runtime communicates variable state to the JavaScript
front end on each step.

The WASM bridge combines the variable table, the current line number, the
execution status, and optional condition information into a single JSON
response per step via \texttt{pseudo\_debug\_step()}.

\paragraph{Condition visualisation.}
When the statement about to execute is a conditional or loop guard, the runtime
additionally captures the condition expression text and its boolean result in a
\texttt{condition\_info\_t} struct stored on the runtime object.  The bridge
layer serialises this alongside the variable state, letting the student see not
only \textit{which} branch was taken but \textit{why} --- the exact condition
text and its evaluated truth value at that step.

% ─────────────────────────────────────────────────────────────────────────────
\section{Transpiler}
\label{sec:transpiler}

\subsection{Code Generation for C, C++, and Pascal}
\label{subsec:codegen}

The transpiler converts a normalised Pseudo++ CST to compilable source code in
C, C++, or Pascal.  Because these three languages differ in significant ways ---
type declarations, operator syntax, string literals, and logical keyword
precedence --- the transpiler abstracts all language-specific choices into a
\texttt{lang\_ops\_t} record populated once at construction time:

\begin{lstlisting}[language=C, basicstyle=\ttfamily\small]
typedef struct {
    const char* assign_op;   /* "=" (C/C++) or ":=" (Pascal) */
    const char* eq_op;       /* "==" or "="  */
    const char* ne_op;       /* "!=" or "<>" */
    const char* or_op;       /* "||" or "or" */
    const char* and_op;      /* "&&" or "and"*/
    const char* not_kw;      /* "!"  or "not " */
    const char* floor_open;  /* "(int)(" or "trunc(" */
    const char* floor_close; /* ")" */
    bool is_pascal;
    bool is_cpp;
} lang_ops_t;
\end{lstlisting}

Transpilation proceeds in two passes:

\paragraph{Pass 1 --- Variable Type Collection.}
\texttt{collect\_vars()} performs a depth-first walk of the CST.  For each
assignment or read statement it records the inferred type of the target
variable in a hash map (\texttt{var\_types}).  The inference rules are:

\begin{itemize}
  \item If the right-hand-side expression contains a string literal, the
        variable is typed as \texttt{string}.
  \item If it contains a floating-point literal (a number with a decimal
        point), it is typed as \texttt{double}.
  \item Otherwise it is typed as \texttt{int}.
\end{itemize}

The first assignment wins for the purpose of code generation; subsequent
assignments may change the value but not the declared type in the target
language.  This is a transpiler-level heuristic, not a language restriction:
the interpreter is fully dynamically typed and allows a variable to hold
values of different types at different points in execution.

\paragraph{Pass 2 --- Code Emission.}
The emission pass walks the CST recursively via \texttt{gen\_stmt()}, using
the \texttt{lang\_ops\_t} record to emit the correct syntax for each construct.
Several language-specific concerns deserve mention:

\begin{itemize}
  \item \textbf{C variable declarations.}  C requires variables to be declared
        before use.  Rather than adding a separate declaration section, the
        transpiler hoists declarations \emph{inline}: the helper
        \texttt{hoist\_vars()} scans each block before emitting its first
        statement and pre-declares any variable first assigned or read in that
        block.  A \texttt{declared\_vars} hash map prevents duplicate
        declarations.  C strings are emitted as \texttt{char name[256]};
        C++ strings use \texttt{std::string}.

  \item \textbf{Pascal variable section.}  For Pascal, variable declarations
        are collected globally during Pass 1 and emitted in a \texttt{var}
        block at the top of the program before the \texttt{begin} keyword.

  \item \textbf{Pascal logical operator precedence.}  In Pascal, \texttt{and}
        and \texttt{or} bind \emph{more tightly} than comparison operators,
        unlike C where \texttt{\&\&} and \texttt{||} bind less tightly.
        To preserve semantics, the transpiler wraps each operand of an
        \texttt{and\_expr} or \texttt{or\_expr} in parentheses when targeting
        Pascal.

  \item \textbf{Pascal string literals.}  String literals are enclosed in
        single quotes in Pascal.  The transpiler re-encodes any embedded
        single quotes using Pascal's doubling convention (\texttt{''}).
\end{itemize}

\subsection{Loop Equivalence Transformations}
\label{subsec:loop-transforms}

The loop equivalence module implements a pedagogically important feature:
given any loop in a Pseudo++ program, it rewrites that loop into an
equivalent form using a different loop construct, grounded in the
interconvertibility established by Böhm and Jacopini \cite{bohm-jacopini}
and the structured programming principles of Dijkstra \cite{dijkstra1968goto}.
This allows students to practise the mechanical equivalence proofs required
at the baccalaureate exam.

The module accepts a source string, a 0-indexed target line identifying the
loop to transform, and the desired output construct (\texttt{FOR},
\texttt{WHILE}, \texttt{DO\_WHILE}, or \texttt{REPEAT}).  It then:

\begin{enumerate}
  \item Runs the source through the formatter and parser to obtain a CST.
  \item Locates the loop node at the specified line via DFS.
  \item Extracts the loop header fields (\texttt{var}, \texttt{start},
        \texttt{end}, \texttt{step}, \texttt{condition}) and body text from
        the CST and the raw source string.
  \item Constructs the target loop form by selecting the appropriate skeleton
        and filling in the extracted fields.
  \item Splices the new loop text back into the original source, preserving
        all preceding and following code unchanged.
\end{enumerate}

Indentation is preserved by detecting the loop's own indentation level from
the source and replicating it in the emitted loop.

% ─────────────────────────────────────────────────────────────────────────────
\section{I/O Abstraction and WebAssembly Integration}
\label{sec:wasm}

\subsection{I/O Abstraction}

The runtime communicates with the outside world exclusively through an abstract
\texttt{io\_t} interface:

\begin{lstlisting}[language=C, basicstyle=\ttfamily\small]
typedef struct {
    const char* (*read)(io_t*);   /* returns next token or NULL */
    void (*write)(io_t*, const char*);
    void (*destroy)(io_t*);
} io_ops_t;

struct io { io_ops_t ops; };
\end{lstlisting}

Two concrete implementations are provided:

\begin{itemize}
  \item \texttt{io\_stdio} --- delegates directly to \texttt{scanf} and
        \texttt{printf}, used by the CLI.
  \item \texttt{io\_buffered} --- maintains an input queue and an output
        queue.  JavaScript pushes input tokens by calling
        \texttt{pseudo\_push\_input()} (a WASM export) and pulls output by
        calling \texttt{pseudo\_pull\_output()}.  When the input queue is empty,
        \texttt{read()} returns \texttt{NULL}, causing the runtime to transition
        to \texttt{EXEC\_NEEDS\_INPUT}.
\end{itemize}

\subsection{WebAssembly Bridge}

The WASM build compiles the entire C library with Emscripten \cite{emscripten},
targeting the WebAssembly instruction format \cite{haas2017wasm}.  The bridge
module (\texttt{wasm/bridge.c}) instantiates a singleton \texttt{runtime\_t}
and a singleton \texttt{io\_buffered}, then exposes the toolchain to
JavaScript via \texttt{EMSCRIPTEN\_KEEPALIVE}-annotated functions, which are
exported by name from the resulting \texttt{.wasm} module.  The exported
surface includes functions for initialisation, source loading, single-step and
run-to-completion execution, input supply, output retrieval, variable
inspection, transpilation, and loop equivalence.

Separating the I/O layer behind an interface means that the \emph{same}
runtime and debugger C code is used in both the CLI and the browser --- only
the I/O backend differs.

% ─────────────────────────────────────────────────────────────────────────────
\section{Web Editor}
\label{sec:editor}

The web editor is a single-page application built around the Monaco Editor
\cite{monaco}, the same editing component that powers Visual Studio Code.
For the \texttt{pseudocode} language, the editor registers:

\begin{itemize}
  \item A \textbf{Monarch tokeniser} \cite{monarch} that classifies tokens
        into keyword, operator, string, number, comment, and identifier
        categories for syntax highlighting.
  \item A \textbf{language configuration} that specifies bracket pairs
        (for auto-closing), the line-comment token, and auto-indentation
        rules.
\end{itemize}

At runtime the JavaScript application loads the compiled WASM module and
communicates with it through the exported C functions.  The interface
provides:

\begin{itemize}
  \item \textbf{Run} --- compiles the editor content, sends it to the WASM
        runtime, and streams output tokens to the output pane.
  \item \textbf{Step / Step Over} --- calls the corresponding WASM step
        function (\texttt{pseudo\_debug\_step} or \texttt{pseudo\_step\_over}),
        which returns a JSON object containing the current line, variable table,
        execution status, and condition information; the front end uses this to
        update the highlighted line and the variable inspector panel.
  \item \textbf{Transpile} --- calls the WASM transpiler with the selected
        target language and displays the result in a read-only Monaco pane.
  \item \textbf{Loop Equivalence} --- identifies the loop under the cursor
        and passes it to the equivalence module.
\end{itemize}

